%!TEX root = main.tex
% ====================================================
\section{Recursive Sample Size Computation}\label{sec:MFMC_mk_recursive}
% ====================================================
   
%%%%%%%%%%%%%%%%%%%%%%%%%%%%%%%%%%%%%%%%%%%%%%%%%%%%%
\subsection{Recursive computation of sample sizes}   
Assume that we have computed integer samples $1 \le m_1 \le \ldots \le m_{j-1}$.
If $m_k = m_{k-1}$ for some $k \in \{ 2, \ldots, j-1 \}$, then model
$k$ is not used and we set $z_k = 0$; otherwise $z_k = 1$.
The remaining computational budget is $b_j = b -  \sum_{k=1}^{j-1} z_k C_k m_k$.

To compute the sample size $m_j^*$,  we consider the problem 
\begin{subequations}\label{eq:Sequential_Optimization}
    \begin{align}
    \label{eq:Sequential_Optimization_obj}
    \min \quad &\sigma_1^2  \sum_{k=j}^K \frac{ \rho_{1,k}^2 - \rho_{1,k+1}^2}{x_k},   \\
       \text{s.t.}\quad & \sum_{k=j}^K C_k x_k \le b_j,       \\
                                & x_j \ge m_{j-1},\quad  x_k \ge x_{k-1}, \quad k=j+1\ldots,K,\label{eq:Sequential_Optimization_obj_c}\\
                                &x_j,\ldots, x_K\in \real.
    \end{align}
\end{subequations}   
The problem  \eqref{eq:Sequential_Optimization} has a similar same structure as 
\eqref{eq:Optimization_sample_size_m_relaxed}, and 
if $j = 1$, $m_0 = 0$, and $b_1 = b$, then \eqref{eq:Sequential_Optimization} is
identical to \eqref{eq:Optimization_sample_size_m_relaxed}.

Because of the monotone correlations assumption \eqref{eq:correlation_assumption},
\eqref{eq:Sequential_Optimization} is a convex optimization problem and has a unique
solution, which we denote by $x_{j|j}^*, \ldots, x_{K|j}^*$.
Theorem~\ref{thm:Sample_size_real} can be used to determine the solution of
\eqref{eq:Sequential_Optimization} with $m_{j-1} = 0$.


\begin{theorem}  \label{thm:Sequential_Optimization_sample_size_real}
   Let the assumption of Theorem~\ref{thm:Sample_size_real} be satisfied and define $\rho_{1,K+1} = 0$.
   \begin{itemize}
   \item[i.] The unique solution of \eqref{eq:Sequential_Optimization} with $m_{j-1} = 0$ is
    \begin{equation}\label{eq:Sequential_Optimization_mjj}
             x_{k|j}^* = \sqrt{\frac{\rho_{1,k}^2 - \rho_{1,k+1}^2}{C_k}} \; 
                          \frac{b_j}{\sum_{l=j}^K \sqrt{C_l (\rho_{1,l}^2 - \rho_{1,l+1}^2)}},
             \quad k = j, \ldots, K,
     \end{equation}
    the cost constraint is active $\sum_{k=1}^K C_kx_k^* = b_j$, 
    and the resulting optimal objective function value is
     \begin{equation}\label{eq:MFMC_variance_optimal-j}
            \sigma_1^2  \sum_{k=j}^K \frac{ \rho_{1,k}^2 - \rho_{1,k+1}^2}{x_{k|j}^*}
           =  \frac{\sigma_1^2}{b_j}\!\left(\sum_{k=j}^K \sqrt{C_k (\rho_{1,k}^2 - \rho_{1,k+1}^2) }\right)^{\!2}.
       \end{equation}
    \item[ii.] If $b_j = b - \sum_{k=1}^{j-1} C_k x_{k|k}^*$, then $x_{j|j}^*$ given by
                  \eqref{eq:Sequential_Optimization_mjj} satisfy $0 < x_{1|1}^* < \ldots < x_{K|K}^*$.
    \item[iii.] If $b_j = b - \sum_{k=1}^{j-1} C_k x_{k|k}^*$, then 
                  then $x_{j|j}^*$ given by \eqref{eq:Sequential_Optimization_mjj} and
                 $x_j^*$ given by \eqref{eq:MFMC_RealValued_Sample_Size} satisfy
                  $x_1^* = x_{1|1}^*,  \ldots, x_K^* = x_{K|K}^*$.
    \end{itemize}
\end{theorem}
\begin{proof}
 \begin{itemize}
   \item[i.] Because  \eqref{eq:Sequential_Optimization} with $m_{j-1} = 0$ has the same structure as 
                \eqref{eq:Optimization_sample_size_m_relaxed}, the first statement follows from
                Theorem~\ref{thm:Sample_size_real}.
  \item[ii.] Define
        \begin{equation}\label{eq:aggregate_cost_variance_weight_S_j}
            S_j = \sum_{k=j}^K\sqrt{C_k (\rho_{1,k}^2 - \rho_{1,k+1}^2) }.
        \end{equation}
        Then 
        \begin{equation}\label{eq:S_j_H_j}
            S_{j-1}=\sqrt{C_{j-1} (\rho_{1,j-1}^2 - \rho_{1,j}^2) }+S_j, \qquad 
            x_{j|j}^* = \sqrt{\frac{\rho_{1,j}^2 - \rho_{1,j-1}^2}{C_j}}\frac{b_j}{S_j},
        \end{equation}
        and
        \begin{equation*}
                     b_j = b_{j-1} - C_{j-1} x_{j-1|j-1}^* 
                         = b_{j-1}\left(1 - \frac{\sqrt{C_{j-1}(\rho_{1,j-1}^2 - \rho_{1,j}^2) }}{S_{j-1}}\right) 
                         = b_{j-1}\frac{S_j}{S_{j-1}}.
        \end{equation*}
        It follows that 
        \begin{equation}\label{eq:b_jS_j}
            \frac{b_j}{S_j} = \frac{b_{j-1}}{S_{j-1}},
        \end{equation}
        and, using assumption \eqref{eq:Sample_size_real_assumptions_b}
        of Theorem~\ref{thm:Sample_size_real},
        \[
             x_{j|j}^* = \sqrt{\frac{\rho_{1,j}^2 - \rho_{1,j+1}^2}{C_j}} \, \frac{b_j}{S_j}
                           = \sqrt{\frac{\rho_{1,j}^2 - \rho_{1,j+1}^2}{C_j}} \,  \frac{b_{j-1}}{S_{j-1}}
                            > \sqrt{\frac{\rho_{1,j-1}^2 - \rho_{1,j}^2}{C_{j-1}}} \, \frac{b_{j-1}}{S_{j-1}}
                           = x_{j-1|j-1}^*.
        \]
        
      \item[iii.]     Recursive application of  \eqref{eq:b_jS_j} gives $b_j = b S_j / S_1$.
             Substituting this expression into \eqref{eq:S_j_H_j} yields
             \[
                        x_{j|j}^*   = \sqrt{\frac{\rho_{1,j}^2 - \rho_{1,j+1}^2}{C_j}} \, \frac{b}{S_1},
             \]
             which coincides with the closed-form MFMC sample size $x_j^*$ in \eqref{eq:MFMC_RealValued_Sample_Size}.  
  \end{itemize}
\end{proof}

Under the conditions of  Theorem~\ref{thm:Sequential_Optimization_sample_size_real},
the components $x_{j|j}^*$ of the solutions $x_{j|j}^*, \ldots, x_{K|j}^*$ of \eqref{eq:Sequential_Optimization_mjj},
$j = 1, \ldots, K-1$,  satisfy $0 < x_1^* = x_{1|1}^* <  \ldots < x_K^* = x_{K|K}^*$.
Thus, the solution $x_{j|j}^*, \ldots, x_{K|j}^*$ of \eqref{eq:Sequential_Optimization_mjj} with $m_{j-1} = 0$ is
also feasible for  \eqref{eq:Sequential_Optimization_mjj} with $m_{j-1} \le x_{j|j}^* = x_j^*$,
and therefore optimal for \eqref{eq:Sequential_Optimization_mjj} with $m_{j-1} \le x_{j|j}^* = x_j^*$.
This leads to the following corollary of Theorem~\ref{thm:Sequential_Optimization_sample_size_real}.

\begin{corollary}  \label{cor:Sequential_Optimization_sample_size_real}
   Let the assumption of Theorem~\ref{thm:Sample_size_real} be satisfied and define $\rho_{1,K+1} = 0$.
   The results of Theorem~\ref{thm:Sequential_Optimization_sample_size_real} remain valid
   if $m_{j-1} = 0$ is replacd by $m_{j-1} \le x_{j|j}^* = x_j^*$.
\end{corollary}



There is a benefit in solving sequences of problems \eqref{eq:Sequential_Optimization} when
one rounds the real-valued sample sizes to integers.

%%%%%%%%%%%%%%%%%%%%%%%%%%%%%%%%%%%%%%%

\subsection{Integer recursive MFMC sample allocation}  
Although the continuous MFMC allocation admits a closed-form optimal solution,
it cannot be implemented directly since sample sizes must be integers.
Solving the resulting integer program exactly becomes increasingly
intractable as the number of fidelity levels $K$ grows.
We therefore construct an integer-valued extension that preserves
the recursive structure of the continuous solution
while maintaining computational efficiency. The idea is to propagate the continuous allocation structure forward through the hierarchy while enforcing integer feasibility at each step. This produces a sequential scheme in which rounding decisions are incorporated into the recursion rather than imposed post hoc, maintaining the monotonicity and budget-consistency properties of the continuous solution. Define the sequential auxiliary real-valued variables
\begin{equation}
\label{eq:MFMC_New_IntegerValued_Sample_Size}
y_1
=
\sqrt{\frac{\Delta_1}{C_1}}\frac{b}{S_1},
\qquad
y_k
=
\sqrt{\frac{\Delta_k}{C_k}}
\frac{b-\sum_{j=1}^{k-1} C_j \lfloor y_j \rfloor}{S_k},
\quad k=2,\ldots,K,
\end{equation}
where $\Delta_k := \rho_{1,k}^2 - \rho_{1,k+1}^2$ and $S_k$ is the aggregate cost–variance weight defined in \eqref{eq:aggregate_cost_variance_weight_S_j}. 



The procedure begins with the continuous allocation $y_1$, from which the integer sample size $\lfloor y_1 \rfloor$ is obtained. The associated cost $C_1 \lfloor y_1 \rfloor$ is deducted from the total budget, and the recursion proceeds to the next fidelity level using the remaining budget. In this manner, each integer allocation $\lfloor y_k \rfloor$ is constructed adaptively with respect to the updated residual budget.

At every iteration, only two quantities are updated: the remaining budget and a cumulative sum over the fidelity hierarchy. Both updates require constant computational effort. Consequently, the per-level computational cost is $\mathcal{O}(1)$, and the overall complexity scales linearly as $\mathcal{O}(K)$ with the number of fidelity levels. Due to this linear scaling and the recursive dependence structure, the scheme avoids any global re-optimization or backtracking. It is therefore particularly well suited for large-scale or high-dimensional fidelity hierarchies, where computational efficiency and scalability are essential.

Following the requirement in \cite{GrGuJuWa:2023}, the total computational budget must satisfy
%
\begin{equation}\label{eq:p_bound}
    b \ge \sum_{k=1}^K C_k,
\end{equation}
%
which guarantees that each fidelity level receives at least one sample, in particular the highest-fidelity model.

Compared with direct flooring of the continuous allocation, the proposed iterative integer scheme achieves improved budget utilization by accounting for cumulative rounding effects across fidelity levels. As a result, the iterative procedure yields integer sample sizes that are no smaller than those obtained by direct rounding. This monotonic dominance property is formalized below.

\begin{corollary}[Monotonic dominance of iterative integer allocation]
\label{thm:MFMC_New_IntegerValued_Cost}
Let $\lfloor y_k \rfloor$ and $\lfloor m_k^* \rfloor$ denote the integer-valued allocations produced by the iterative scheme and by direct flooring of the continuous optimal allocation, respectively. Under the budget condition \eqref{eq:p_bound}, the allocations satisfy
\[
\lfloor m_k^* \rfloor \le \lfloor y_k \rfloor,
\qquad k=1,\ldots,K.
\]
\end{corollary}
%

\begin{proof}
We first show that the total cost of the iterative scheme does not exceed the prescribed budget. Define the cumulative integer cost up to level $k$ as
\[
T_k = \sum_{j=1}^k C_j\left\lfloor y_j \right\rfloor.
\]
Since $\lfloor y_j \rfloor \le y_j$, we claim, and prove by induction, that for each $k = 1, \ldots, K$,
\begin{equation}\label{eq:Tk_bound}
T_k \le \frac{b}{S_1}\sum_{j=1}^k \sqrt{C_j \Delta_j}.
\end{equation}
where $S_k$ is the aggregate cost–variance weight defined in \eqref{eq:aggregate_cost_variance_weight_S_j}. 
Inequality \eqref{eq:Tk_bound} bounds the cumulative integer cost at level $k$ by a proportional share of the total budget, scaled by $S_1$.


The base case $k=1$ holds for \eqref{eq:Tk_bound} immediately,
\[
T_1=C_1 \left\lfloor y_1 \right\rfloor \le C_1y_1 = \frac{b}{S_1}\sqrt{C_1\Delta_1}.
\]
Assume \eqref{eq:Tk_bound} holds for \(k-1\). By definition of \(y_k\),
%
\[
y_k = \sqrt{\frac{\Delta_k}{C_k}}\frac{b - T_{k-1}}{S_k},
\]
%
and hence
%
\[
C_k \left\lfloor y_k \right\rfloor \le C_k y_k  = \sqrt{C_k\Delta_k}\frac{b-T_{k-1}}{S_k}.
\]
%
Using the inductive hypothesis and simplifying yields
\begin{align*}
    T_k &= T_{k-1}+C_k\left\lfloor y_k \right\rfloor \\
    &\le T_{k-1} + \sqrt{C_k\Delta_k}\frac{b-T_{k-1}}{S_k}
    =T_{k-1}\left(1-\frac{\sqrt{C_k\Delta_k}}{S_k}\right) + b\frac{\sqrt{C_k\Delta_k}}{S_k}
    \le \frac{b}{S_1}\sum_{j=1}^{k-1} \sqrt{C_j\Delta_j}\frac{S_{k+1}}{S_k}+b\frac{\sqrt{C_k\Delta_k}}{S_k}\\
    % =\frac{p}{S_1}\cdot \frac{\sum_{j=1}^{k-1} \sqrt{C_j\Delta_j}S_{k+1}+S_1\sqrt{C_k\Delta_k}}{S_k}
    &=\frac{b}{S_1}\cdot \frac{\sum_{j=1}^{k-1} \sqrt{C_j\Delta_j}S_{k+1}+\sqrt{C_k\Delta_k}\left(\sum_{j=1}^{k-1}\sqrt{C_j\Delta_j}+S_k\right)}{S_k}
    %=\frac{p}{S}\frac{\sum_{j=k}^K\sqrt{C_j\Delta_j}\sum_{j=1}^k\sqrt{C_j\Delta_j}}{\sum_{j=k}^K\sqrt{C_j\Delta_j}}
    =\frac{b}{S_1}\sum_{j=1}^k\sqrt{C_j\Delta_j}.
\end{align*}
%
Thus, inequality \eqref{eq:Tk_bound} holds for all $k$. 

To establish the lower bound in \eqref{eq:Iterative_integer_sample_size_cost_bound}, we compare the auxiliary sequences $y_k$ and $m_k^*$. From \eqref{eq:Tk_bound},
%
\[
y_k = \sqrt{\frac{\Delta_k}{C_k}}\frac{b - T_{k-1}}{S_k} \ge \sqrt{\frac{\Delta_k}{C_k}}\frac{b-\frac{b}{S}\sum_{j=1}^{k-1}\sqrt{C_j\Delta_j}}{S_k} = \sqrt{\frac{\Delta_k}{C_k}}\frac{b}{S_1}=m_k^*, \qquad k \ge 1.
\]
% 
Monotonicity of the floor function implies \(\lfloor y_k\rfloor\ge\lfloor m_k^*\rfloor\) for every \(k\). 

\end{proof}


Using Corollary~\ref{thm:MFMC_New_IntegerValued_Cost}, we show that the iterative integer-valued allocation achieves improved budget utilization by compensating for the cumulative rounding effects introduced by the floor operation. In contrast to direct flooring of the continuous allocation, which may leave unused budget and thereby inflate the variance, the iterative scheme redistributes the residual budget across subsequent fidelity levels. As a result, it mitigates variance inflation caused by budget under-utilization and yields superior variance performance and cost efficiency. These properties are quantified in Lemma~\ref{thm:MFMC_New_IntegerValued_Variance}.

\begin{lemma}[Variance improvement and budget efficiency of the iterative integer allocation]
\label{thm:MFMC_New_IntegerValued_Variance}
Under the computational budget constraint \eqref{eq:p_bound}, the iterative integer-valued allocation $\{\lfloor y_k \rfloor\}_{k=1}^K$ satisfies:

\medskip
\noindent
\textbf{(i) Budget efficiency.}
\begin{equation}\label{eq:Iterative_integer_sample_size_cost_bound}
\sum_{k=1}^K C_k \left\lfloor m_k^* \right\rfloor
\;\le\;
\sum_{k=1}^K C_k \left\lfloor y_k \right\rfloor
\;\le\;
b,
\end{equation}
that is, the iterative allocation never exceeds the prescribed budget and always utilizes at least as much budget as direct flooring.

\medskip
\noindent
\textbf{(ii) Variance improvement.}
\begin{equation}\label{eq:Iterative_Integer_Variance_Bound}
\frac{1}{b}
\left(
\sum_{k=1}^K \sqrt{C_k \Delta_k}
\right)^2=:\sum_{k=1}^K \frac{\Delta_k}{m_k^*}
\;\le\;
\sum_{k=1}^K \frac{\Delta_k}{\left\lfloor y_k \right\rfloor}
\;\le\;
\sum_{k=1}^K \frac{\Delta_k}{\left\lfloor m_k^* \right\rfloor},
\end{equation}

In particular, the iterative integer allocation strictly reduces the variance inflation caused by direct flooring.


\end{lemma}







\begin{proof}
\medskip
\noindent
\mbox{}\\
{\it Cost bounds.}
From the construction of the iterative scheme (Lemma~\ref{thm:MFMC_New_IntegerValued_Cost}), we have
\[
\left\lfloor y_k \right\rfloor \;\ge\; \left\lfloor m_k^* \right\rfloor,
\qquad k=1,\dots,K.
\]
Multiplying by $C_k$ and summing over $k$ immediately yields the lower bound in \eqref{eq:Iterative_integer_sample_size_cost_bound}. By construction of the iterative allocation, the cumulative cost after the final level satisfies
\begin{equation}\label{eq:MFMC_iterative_total_cost}
T_K
:=
\sum_{k=1}^K C_k \left\lfloor y_k \right\rfloor
\le b,
\end{equation}
which establishes the upper bound in \eqref{eq:Iterative_integer_sample_size_cost_bound}.


\medskip
\noindent
\textit{Equality conditions.}
The upper bound is attained if and only if no budget remains unused, i.e.,
\[
\sum_{k=1}^K C_k
\left(y_k - \left\lfloor y_k \right\rfloor\right) = 0.
\]
Since each fractional part $y_k - \lfloor y_k \rfloor $ lies in $[0,1)$, this equality condition occurs precisely when every $y_k$ is integer.  
The lower bound is attained if and only if \(\lfloor y_k \rfloor = \lfloor m_k^* \rfloor\) for all \(k\), equivalently when both real-valued allocations \(y_k\) and \(m_k^*\) lie in the same integer interval, i.e. $
m_k^* \le y_k < \left\lfloor m_k^* \right\rfloor + 1$, for all $k=1,\dots,K.$

\bigskip
\paragraph{Variance bounds}
We first prove the upper bound in \eqref{eq:Iterative_Integer_Variance_Bound}.  
Since $\lfloor y_k \rfloor \ge \lfloor m_k^* \rfloor$ from Lemma~\ref{thm:MFMC_New_IntegerValued_Cost} and the function
$x \mapsto \Delta_k/x$ is strictly decreasing for $x>0$, we immediately obtain
%
\[
\sum_{k=1}^K
\frac{\Delta_k}{\left\lfloor y_k \right\rfloor}
\le
\sum_{k=1}^K
\frac{\Delta_k}{\left\lfloor m_k^* \right\rfloor}.
\]

To establish the lower bound, we apply the Cauchy--Schwarz inequality:
\[
\left(\sum_{k=1}^K \sqrt{C_k \Delta_k}\right)^2
\le
\left(\sum_{k=1}^K C_k \left\lfloor y_k \right\rfloor\right)
\left(\sum_{k=1}^K \frac{\Delta_k}{\left\lfloor y_k \right\rfloor}\right).
\]
Using \eqref{eq:MFMC_iterative_total_cost}, we have $\sum_{k=1}^K C_k \left\lfloor y_k \right\rfloor \le b$, and therefore
\[
\frac{1}{b}
\left(\sum_{k=1}^K \sqrt{C_k \Delta_k}\right)^2
\le
\sum_{k=1}^K
\frac{\Delta_k}{\left\lfloor y_k \right\rfloor}.
\]

\medskip
\noindent
\textit{Equality conditions.}
Equality in the Cauchy--Schwarz step holds if and only if there exists
$\lambda>0$ such that
\[
\left\lfloor y_k \right\rfloor
=
\lambda \sqrt{\frac{\Delta_k}{C_k}},
\qquad k=1,\dots,K,
\]
which indicates $\lfloor y_k \rfloor= m_k^*$.  
Thus equality in \eqref{eq:Iterative_Integer_Variance_Bound} holds if and only if the iterative scheme reproduces the continuous solution exactly.






\end{proof}



% ------------------
% Next consider the variance
% \begin{align*}
%     f_{\text{act}}(\overline{N_k})&=\sum_{i=1}^{K}\frac{\Delta_k}{\overline{N_k}}=\sum_{i=1}^{k-1}\frac{\Delta_i}{\overline{N_i}}+\frac{\Delta_k}{\overline{N_k}}+\sum_{i=k+1}^{K}\frac{\Delta_i}{\overline{N_i}}\\
%     &\in \left[\sum_{i=1}^{k-1}\frac{\Delta_i}{\overline{N_i}}+\frac{\Delta_k}{\overline{N_k}}+\sum_{i=k+1}^K\frac{\Delta_{i}}{N_i^*},\; \sum_{i=1}^{k-1}\frac{\Delta_i}{\overline{N_i}}+\frac{\Delta_k}{\overline{N_k}}+\sum_{i=k+1}^K\frac{\Delta_{i}}{N_i^*-1}\right)=[f_1,f_2)\\
%     f_1&=\sum_{i=1}^{k-1}\frac{\Delta_i}{\overline{N_i}}+\frac{\Delta_k}{\overline{N_k}}+\sum_{i=k+1}^K\frac{\Delta_{i}}{N_i^*}=\sum_{i=1}^{k-1}\frac{\Delta_i}{\overline{N_i}}+\frac{\Delta_k}{\overline{N_k}}+\sum_{i=k+1}^K\sqrt{C_i\Delta_i}\frac{\sum_{j=i}^K\sqrt{C_j\Delta_j}}{p-\sum_{j=1}^{i-1}C_j\overline{N_j}}\\
%     f_2&=\sum_{i=1}^{k-1}\frac{\Delta_i}{\overline{N_i}}+\frac{\Delta_k}{\overline{N_k}}+\sum_{i=k+1}^K \ldots\\
% \end{align*}
% \begin{align*}
%     \frac{d f_1}{d \overline{N_k}}&=-\frac{\Delta_k}{\overline{N_k}^2}+C_k\sum_{i=k+1}^K\sqrt{C_i\Delta_i}\frac{\sum_{j=i}^K\sqrt{C_j\Delta_j}}{\left(p-\sum_{j=1}^{i-1}C_j\overline{N_j}\right)^2}\\
%     \frac{d^2 f_1}{d^2 \overline{N_k}}&=\frac{2\Delta_k}{\overline{N_k}^3}+2C_k^2\sum_{i=k+1}^K\sqrt{C_i\Delta_i}\frac{\sum_{j=i}^K\sqrt{C_j\Delta_j}}{\left(p-\sum_{j=1}^{i-1}C_j\overline{N_j}\right)^3}
% \end{align*}
% Note that $\frac{d^2 f_1}{d^2 \overline{N_k}}>0$ whenever $p>\sum_{j=1}^{i-1}C_j\overline{N_j}$. this means $f_1$ is convex in $\overline{N_k}$.